% Beide Ausgabearten führen dieselben Hauptdeklarationen in gleicher Reihenfolge aus.
\CSBPartDef
\CSBProof{%
\Needspace{14\baselineskip}
}

\CSBPartSubset
\CSBProof{%
\begin{tabproof}
  \proofstep{1}{F\colon A\inj B}{\rA}
  \proofstep{2}{G\colon B\inj A}{\rA}
  \proofstep{3}{x\in\CBPart{F}{G}}{\rA}
  \proofstep{1,2,3}{x\in A}{%
    \FormulaRefAuto{CantorBernsteinPartDef}[def]{1,2,3}}
  \proofstep{1,2}{\CBPart{F}{G}\subseteq A}{%
    \FormulaRefAuto{SubsetIntroductionRule}[thm]{[3]\vdots 4}}
\end{tabproof}
}
\Needspace{14\baselineskip}

\CSBPartContainsSeed
\CSBProof{%
\begin{tabproofwide}
  \proofstepwidestar[1]{F\colon A\inj B}{\rA}
  \proofstepwidestar[2]{G\colon B\inj A}{\rA}
  \proofstepwidestar[3]{x\in A\setminus G[B]}{\rA}
  \proofstepwidestar[4]{%
    (A\setminus G[B])\subseteq X\land G[F[X]]\subseteq X}{\rA}
  \proofstepwidestar[4]{A\setminus G[B]\subseteq X}{\rAEa{4}}
  \proofstepwidestar[3,4]{x\in X}{%
    \FormulaRefAuto{A\subseteq B\dsep x\in A\vdash x\in B}{5,3}}
  \proofstepwidestar[3]{%
    ((A\setminus G[B])\subseteq X\land G[F[X]]\subseteq X)
      \rightarrow x\in X}{\rRI{4,6}}
  \proofstepwidestar[3]{%
    X\in\powerset(A)\rightarrow\bigl(
      ((A\setminus G[B])\subseteq X\land G[F[X]]\subseteq X)
        \rightarrow x\in X\bigr)}{%
    \FormulaRefAuto{Q \vdash P \rightarrow Q}[thm]{7}}
  \proofstepwidestar[3]{%
    \forall X\in\powerset(A)\,\bigl(
      ((A\setminus G[B])\subseteq X\land G[F[X]]\subseteq X)
      \rightarrow x\in X\bigr)}{\rUI{8}}
  \proofstepwidestar[3]{x\in A}{%
    \FormulaRefAuto{x\in A\setminus B\vdash x\in A}{3}}
  \proofstepwidestar[1,2,3]{x\in\CBPart{F}{G}}{%
    \FormulaRefAuto{CantorBernsteinPartDef}[def]{1,2,\rAI{10,9}}}
  \proofstepwidestar[1,2]{A\setminus G[B]\subseteq\CBPart{F}{G}}{%
    \FormulaRefAuto{SubsetIntroductionRule}[thm]{[3]\vdots 11}}
\end{tabproofwide}
}
\Needspace{14\baselineskip}

\CSBPartMinimal
\CSBProof{%
\begin{tabproofwide}
  \proofstepwidestar[1]{F\colon A\inj B}{\rA}
  \proofstepwidestar[2]{G\colon B\inj A}{\rA}
  \proofstepwidestar[3]{X\subseteq A}{\rA}
  \proofstepwidestar[4]{A\setminus G[B]\subseteq X}{\rA}
  \proofstepwidestar[5]{G[F[X]]\subseteq X}{\rA}
  \proofstepwidestar[6]{x\in\CBPart{F}{G}}{\rA}
  \proofstepwidestar[3]{X\in\powerset(A)}{%
    \FormulaRefAuto{X\in\powerset(A)\eqvdash X\subseteq A}{3}}
  \proofstepwidestar[4,5]{%
    (A\setminus G[B])\subseteq X\land G[F[X]]\subseteq X}{\rAI{4,5}}
  \proofstepwidestar[1,2]{%
    \CBPart{F}{G}=\bigl\{z\in A\bigm|
      \forall Y\in\powerset(A)\,\bigl(
        ((A\setminus G[B])\subseteq Y\land G[F[Y]]\subseteq Y)
          \rightarrow z\in Y\bigr)\bigr\}}{%
    \FormulaRefAuto{CantorBernsteinPartDef}[def]{1,2}}
  \proofstepwidestar[1,2,6]{%
    \forall Y\in\powerset(A)\,\bigl(
      ((A\setminus G[B])\subseteq Y\land G[F[Y]]\subseteq Y)
        \rightarrow x\in Y\bigr)}{%
    \FormulaRefAuto{x \in A\dsep A=\{x \in B \mid P(x)\}
      \vdash P(x)}[thm]{6,9}}
  \proofstepwidestar[1,2,6]{%
    X\in\powerset(A)\rightarrow\bigl(
      ((A\setminus G[B])\subseteq X\land G[F[X]]\subseteq X)
        \rightarrow x\in X\bigr)}{\rUE{10}}
  \proofstepwidestar[1,2,3,6]{%
    ((A\setminus G[B])\subseteq X\land G[F[X]]\subseteq X)
      \rightarrow x\in X}{\rRE{11,7}}
  \proofstepwidestar[1,2,3,4,5,6]{x\in X}{\rRE{12,8}}
  \proofstepwidestar[1,2,3,4,5]{\CBPart{F}{G}\subseteq X}{%
    \FormulaRefAuto{SubsetIntroductionRule}[thm]{[6]\vdots 13}}
\end{tabproofwide}
}
\Needspace{24\baselineskip}

\CSBPartClosed
\CSBProof{%
\begin{tabproofsplitwide}
  \Needspace{8\baselineskip}
  \hypertarget{csb.helper.CantorBernsteinPartClosedUniversal}{}
  \proofpartwideindR[Bildpunkte liegen in jeder zulässigen Teilmenge]{%
    \begin{aligned}[t]
      &F\colon A\inj B\dsep G\colon B\inj A\dsep
        y\in G[F[\CBPart{F}{G}]]\vdash{}\\[-2pt]
      &\forall X\in\powerset(A)\,\bigl(
        ((A\setminus G[B])\subseteq X\land G[F[X]]\subseteq X)
        \rightarrow y\in X\bigr)
    \end{aligned}%
  }{%
    F\colon A\inj B\dsep G\colon B\inj A\dsep
    y\in G[F[\CBPart{F}{G}]]\vdash
    \forall X\in\powerset(A)\,\bigl(
      ((A\setminus G[B])\subseteq X\land G[F[X]]\subseteq X)
      \rightarrow y\in X\bigr)
  }[CantorBernsteinPartClosedUniversal]
    \proofstepwidestar[1]{F\colon A\inj B}{\rA}
    \proofstepwidestar[2]{G\colon B\inj A}{\rA}
    \proofstepwidestar[3]{y\in G[F[\CBPart{F}{G}]]}{\rA}
    \proofstepwidestar[4]{X\in\powerset(A)}{\rA}
    \proofstepwidestar[5]{%
      (A\setminus G[B])\subseteq X\land G[F[X]]\subseteq X}{\rA}
    \proofstepwidestar[4]{X\subseteq A}{%
      \FormulaRefAuto{X\in\powerset(A)\eqvdash X\subseteq A}{4}}
    \proofstepwidestar[5]{A\setminus G[B]\subseteq X}{\rAEa{5}}
    \proofstepwidestar[5]{G[F[X]]\subseteq X}{\rAEb{5}}
    \proofstepwidestar[1,2,4,5]{\CBPart{F}{G}\subseteq X}{%
      \FormulaRefAuto{CantorBernsteinPartMinimal}[thm]{1,2,6,7,8}}
    \proofstepwidestar[1]{F\colon A\to B}{%
      \FormulaRefAuto{Injektive Funktion}[def]{1}}
    \proofstepwidestar[2]{G\colon B\to A}{%
      \FormulaRefAuto{Injektive Funktion}[def]{2}}
    \proofstepwidestar[1,2,4,5]{F[\CBPart{F}{G}]\subseteq F[X]}{%
      \FormulaRefAuto{F\colon M\to N\dsep A\subseteq B\dsep B\subseteq M
        \vdash F[A]\subseteq F[B]}[thm]{10,9,6}}
    \proofstepwidestar[1,4]{F[X]\subseteq B}{%
      \FormulaRefAuto{C\subseteq A\dsep F\colon A\to B
        \vdash F[C]\subseteq B}[thm]{6,10}}
    \proofstepwidestar[1,2,4,5]{%
      G[F[\CBPart{F}{G}]]\subseteq G[F[X]]}{%
      \FormulaRefAuto{F\colon M\to N\dsep A\subseteq B\dsep B\subseteq M
        \vdash F[A]\subseteq F[B]}[thm]{11,12,13}}
    \proofstepwidestar[1,2,3,4,5]{y\in G[F[X]]}{%
      \FormulaRefAuto{A\subseteq B\dsep x\in A\vdash x\in B}[thm]{14,3}}
    \proofstepwidestar[1,2,3,4,5]{y\in X}{%
      \FormulaRefAuto{A\subseteq B\dsep x\in A\vdash x\in B}[thm]{8,15}}
    \proofstepwidestar[1,2,3,4]{%
      ((A\setminus G[B])\subseteq X\land G[F[X]]\subseteq X)
        \rightarrow y\in X}{\rRI{5,16}}
    \proofstepwidestar[1,2,3]{%
      \forall X\in\powerset(A)\,\bigl(
        ((A\setminus G[B])\subseteq X\land G[F[X]]\subseteq X)
        \rightarrow y\in X\bigr)}{\rUI{\rRI{4,17}}}
  \closeproofpartwideindR

  \Needspace{8\baselineskip}

  \proofpartwide{Zugehörigkeit zu \(A\) und Schluss}
    \proofstepwidestar[1]{F\colon A\inj B}{\rA}
    \proofstepwidestar[2]{G\colon B\inj A}{\rA}
    \proofstepwidestar[3]{y\in G[F[\CBPart{F}{G}]]}{\rA}
    \proofstepwidestar[1,2,3]{%
      \forall X\in\powerset(A)\,\bigl(
        ((A\setminus G[B])\subseteq X\land G[F[X]]\subseteq X)
        \rightarrow y\in X\bigr)}{%
      \FormulaRefAuto{CantorBernsteinPartClosedUniversal}[thm]{1,2,3}}
    \proofstepwidestar[1]{F\colon A\to B}{%
      \FormulaRefAuto{Injektive Funktion}[def]{1}}
    \proofstepwidestar[2]{G\colon B\to A}{%
      \FormulaRefAuto{Injektive Funktion}[def]{2}}
    \proofstepwidestar[1,2]{\CBPart{F}{G}\subseteq A}{%
      \FormulaRefAuto{CantorBernsteinPartSubset}[thm]{1,2}}
    \proofstepwidestar[1,2]{F[\CBPart{F}{G}]\subseteq B}{%
      \FormulaRefAuto{C\subseteq A\dsep F\colon A\to B
        \vdash F[C]\subseteq B}[thm]{7,5}}
    \proofstepwidestar[1,2]{G[F[\CBPart{F}{G}]]\subseteq A}{%
      \FormulaRefAuto{C\subseteq A\dsep F\colon A\to B
        \vdash F[C]\subseteq B}[thm]{8,6}}
    \proofstepwidestar[1,2,3]{y\in A}{%
      \FormulaRefAuto{A\subseteq B\dsep x\in A\vdash x\in B}[thm]{9,3}}
    \proofstepwidestar[1,2,3]{y\in\CBPart{F}{G}}{%
      \FormulaRefAuto{CantorBernsteinPartDef}[def]{1,2,\rAI{10,4}}}
    \proofstepwidestar[1,2]{%
      G[F[\CBPart{F}{G}]]\subseteq\CBPart{F}{G}}{%
      \FormulaRefAuto{SubsetIntroductionRule}[thm]{[3]\vdots 11}}
  \closeproofpartwide
\end{tabproofsplitwide}
}
\Needspace{14\baselineskip}

\CSBPartFixedPoint
\CSBProof{%
\begin{tabproofsplitwide}
\Needspace{8\baselineskip}
\proofpartwide{Erste Inklusion}
\proofstepwidestar[1]{F\colon A\inj B}{\rA}
  \proofstepwidestar[2]{G\colon B\inj A}{\rA}
  \proofstepwidestar[1,2]{A\setminus G[B]\subseteq\CBPart{F}{G}}{%
    \FormulaRefAuto{CantorBernsteinPartContainsSeed}[thm]{1,2}}
  \proofstepwidestar[1,2]{%
    G[F[\CBPart{F}{G}]]\subseteq\CBPart{F}{G}}{%
    \FormulaRefAuto{CantorBernsteinPartClosed}[thm]{1,2}}
  \proofstepwidestar[1,2]{%
    (A\setminus G[B])\cup G[F[\CBPart{F}{G}]]
      \subseteq\CBPart{F}{G}}{%
    \FormulaRefAuto{UnionSubsetCommonSuperset}[thm]{3,4}}

\closeproofpartwide

\Needspace{8\baselineskip}

\proofpartwide{Zweite Inklusion}
  \setcounter{proofstepnr}{5}% Fortlaufende Schritte dieses Beweises.
\proofstepwidestar[]{%
    A\setminus G[B]\subseteq
      (A\setminus G[B])\cup G[F[\CBPart{F}{G}]]}{%
    \FormulaRefAuto{A\subseteq A\cup B}[thm]}
  \proofstepwidestar[1,2]{\CBPart{F}{G}\subseteq A}{%
    \FormulaRefAuto{CantorBernsteinPartSubset}[thm]{1,2}}
  \proofstepwidestar[1,2]{%
    (A\setminus G[B])\cup G[F[\CBPart{F}{G}]]\subseteq A}{%
    \FormulaRefAuto{A\subseteq B\dsep B\subseteq C\vdash A\subseteq C}[thm]{5,7}}
  \proofstepwidestar[1]{F\colon A\to B}{%
    \FormulaRefAuto{Injektive Funktion}[def]{1}}
  \proofstepwidestar[2]{G\colon B\to A}{%
    \FormulaRefAuto{Injektive Funktion}[def]{2}}
  \proofstepwidestar[1,2]{%
    F[(A\setminus G[B])\cup G[F[\CBPart{F}{G}]]]
      \subseteq F[\CBPart{F}{G}]}{%
    \FormulaRefAuto{F\colon M\to N\dsep A\subseteq B\dsep B\subseteq M
      \vdash F[A]\subseteq F[B]}[thm]{9,5,7}}
  \proofstepwidestar[1,2]{F[\CBPart{F}{G}]\subseteq B}{%
    \FormulaRefAuto{C\subseteq A\dsep F\colon A\to B
      \vdash F[C]\subseteq B}[thm]{7,9}}
  \proofstepwidestar[1,2]{%
    G[F[(A\setminus G[B])\cup G[F[\CBPart{F}{G}]]]]
      \subseteq G[F[\CBPart{F}{G}]]}{%
    \FormulaRefAuto{F\colon M\to N\dsep A\subseteq B\dsep B\subseteq M
      \vdash F[A]\subseteq F[B]}[thm]{10,11,12}}
  \proofstepwidestar[]{%
    G[F[\CBPart{F}{G}]]\subseteq
      (A\setminus G[B])\cup G[F[\CBPart{F}{G}]]}{%
    \FormulaRefAuto{B\subseteq A\cup B}[thm]}
  \proofstepwidestar[1,2]{%
    G[F[(A\setminus G[B])\cup G[F[\CBPart{F}{G}]]]]
      \subseteq (A\setminus G[B])\cup G[F[\CBPart{F}{G}]]}{%
    \FormulaRefAuto{A\subseteq B\dsep B\subseteq C\vdash A\subseteq C}[thm]{13,14}}
  \proofstepwidestar[1,2]{%
    \CBPart{F}{G}\subseteq
      (A\setminus G[B])\cup G[F[\CBPart{F}{G}]]}{%
    \FormulaRefAuto{CantorBernsteinPartMinimal}[thm]{1,2,8,6,15}}

\closeproofpartwide

\Needspace{8\baselineskip}

\proofpartwide{Fixpunktgleichheit}
  \setcounter{proofstepnr}{16}% Fortlaufende Schritte dieses Beweises.
\proofstepwidestar[1,2]{%
    \CBPart{F}{G}=(A\setminus G[B])\cup G[F[\CBPart{F}{G}]]}{%
    \FormulaRefAuto{A\subseteq B\dsep B\subseteq A\vdash A=B}[thm]{16,5}}

\closeproofpartwide
\end{tabproofsplitwide}
}
\Needspace{14\baselineskip}

\CSBComplementIdentity
\CSBProof{%
\begin{tabproofsplitwide}
  \Needspace{8\baselineskip}
  \hypertarget{csb.helper.CantorBernsteinComplementForward}{}
  \proofpartwideindR[Teilmengenrichtung \(\subseteq\)]{%
    F\colon A\inj B\dsep G\colon B\inj A\vdash
    G[B]\setminus G[F[\CBPart{F}{G}]]
      \subseteq A\setminus\CBPart{F}{G}
  }{%
    F\colon A\inj B\dsep G\colon B\inj A\vdash
    G[B]\setminus G[F[\CBPart{F}{G}]]
      \subseteq A\setminus\CBPart{F}{G}
  }[CantorBernsteinComplementForward]
    \proofstepwidestar[1]{F\colon A\inj B}{\rA}
    \proofstepwidestar[2]{G\colon B\inj A}{\rA}
    \proofstepwidestar[3]{y\in G[B]\setminus G[F[\CBPart{F}{G}]]}{\rA}
    \proofstepwidestar[3]{y\in G[B]}{%
      \FormulaRefAuto{x\in A\setminus B\vdash x\in A}[thm]{3}}
    \proofstepwidestar[3]{y\notin G[F[\CBPart{F}{G}]]}{%
      \FormulaRefAuto{x\in A\setminus B\vdash x\notin B}[thm]{3}}
    \proofstepwidestar[2]{G\colon B\to A}{%
      \FormulaRefAuto{Injektive Funktion}[def]{2}}
    \proofstepwidestar[2]{G[B]\subseteq A}{%
      \FormulaRefAuto{F\colon A\to B\vdash F[A]\subseteq B}[thm]{6}}
    \proofstepwidestar[2,3]{y\in A}{%
      \FormulaRefAuto{A\subseteq B\dsep x\in A\vdash x\in B}[thm]{7,4}}
  \proofstepwidestar[2,3]{( y \notin A \setminus { G [ B ] } \leftrightarrow y \in { G [ B ] } )}{%
      \FormulaRefAuto{RelCompNotMembership}[thm]{8}}
  \proofstepwidestar[2,3]{y\notin A\setminus G[B]}{%
      \FormulaRefAuto{P \leftrightarrow Q, Q \vdash P}[thm]{%
        9,4}}
    \proofstepwidestar[1,2]{%
      \CBPart{F}{G}=(A\setminus G[B])\cup G[F[\CBPart{F}{G}]]}{%
      \FormulaRefAuto{CantorBernsteinPartFixedPoint}[thm]{1,2}}
  \proofstepwidestar[1,2,3]{%
    y\notin(A\setminus G[B])\cup G[F[\CBPart{F}{G}]]}{%
      \FormulaRefAuto{NotInUnion}[thm]{10,5}}
    \proofstepwidestar[1,2,3]{y\notin\CBPart{F}{G}}{\rIE{11,12}}
    \proofstepwidestar[1,2,3]{y\in A\setminus\CBPart{F}{G}}{%
      \FormulaRefAuto{x\in A\setminus B\eqvdash x\in A\land x\notin B}[thm]{%
        \rAI{8,13}}}
    \proofstepwidestar[1,2]{%
      G[B]\setminus G[F[\CBPart{F}{G}]]
        \subseteq A\setminus\CBPart{F}{G}}{%
      \FormulaRefAuto{SubsetIntroductionRule}[thm]{[3]\vdots 14}}
  \closeproofpartwideindR

  \Needspace{8\baselineskip}

  \hypertarget{csb.helper.CantorBernsteinComplementBackward}{}
  \proofpartwideindR[Teilmengenrichtung \(\supseteq\)]{%
    F\colon A\inj B\dsep G\colon B\inj A\vdash
    A\setminus\CBPart{F}{G}
      \subseteq G[B]\setminus G[F[\CBPart{F}{G}]]
  }{%
    F\colon A\inj B\dsep G\colon B\inj A\vdash
    A\setminus\CBPart{F}{G}
      \subseteq G[B]\setminus G[F[\CBPart{F}{G}]]
  }[CantorBernsteinComplementBackward]
    \proofstepwidestar[1]{F\colon A\inj B}{\rA}
    \proofstepwidestar[2]{G\colon B\inj A}{\rA}
    \proofstepwidestar[3]{y\in A\setminus\CBPart{F}{G}}{\rA}
    \proofstepwidestar[3]{y\in A}{%
      \FormulaRefAuto{x\in A\setminus B\vdash x\in A}[thm]{3}}
    \proofstepwidestar[3]{y\notin\CBPart{F}{G}}{%
      \FormulaRefAuto{x\in A\setminus B\vdash x\notin B}[thm]{3}}
    \proofstepwidestar[1,2]{%
      \CBPart{F}{G}=(A\setminus G[B])\cup G[F[\CBPart{F}{G}]]}{%
      \FormulaRefAuto{CantorBernsteinPartFixedPoint}[thm]{1,2}}
    \proofstepwidestar[1,2,3]{%
      y\notin(A\setminus G[B])\cup G[F[\CBPart{F}{G}]]}{\rIE{6,5}}
  \proofstepwidestar[1,2,3]{y \notin { ( A \setminus G [ B ] ) } \land y \notin { G [ F [ \CBPart F G ] ] }}{%
      \FormulaRefAuto{OutsideUnionComponents}[thm]{7}}
  \proofstepwidestar[1,2,3]{y\notin A\setminus G[B]}{%
      \rAEa{8}}
  \proofstepwidestar[1,2,3]{y \notin { ( A \setminus G [ B ] ) } \land y \notin { G [ F [ \CBPart F G ] ] }}{%
      \FormulaRefAuto{OutsideUnionComponents}[thm]{7}}
  \proofstepwidestar[1,2,3]{y\notin G[F[\CBPart{F}{G}]]}{%
      \rAEb{10}}
    \proofstepwidestar[12]{y\notin G[B]}{\rA}
    \proofstepwidestar[3,12]{y\in A\setminus G[B]}{%
      \FormulaRefAuto{x\in A\setminus B\eqvdash x\in A\land x\notin B}[thm]{%
        \rAI{4,12}}}
    \proofstepwidestar[1,2,3,12]{\bot}{\rBI{13,9}}
    \proofstepwidestar[1,2,3]{\neg(y\notin G[B])}{\rCI{12,14}}
    \proofstepwidestar[1,2,3]{y\in G[B]}{\FormulaRefAuto{rDN}[thm]{15}}
    \proofstepwidestar[1,2,3]{%
      y\in G[B]\setminus G[F[\CBPart{F}{G}]]}{%
      \FormulaRefAuto{x\in A\setminus B\eqvdash x\in A\land x\notin B}[thm]{%
        \rAI{16,11}}}
    \proofstepwidestar[1,2]{%
      A\setminus\CBPart{F}{G}
        \subseteq G[B]\setminus G[F[\CBPart{F}{G}]]}{%
      \FormulaRefAuto{SubsetIntroductionRule}[thm]{[3]\vdots 17}}
  \closeproofpartwideindR

  \Needspace{8\baselineskip}

  \proofpartwide{Schluss}
    \proofstepwidestar[1]{F\colon A\inj B}{\rA}
    \proofstepwidestar[2]{G\colon B\inj A}{\rA}
    \proofstepwidestar[1,2]{%
      G[B]\setminus G[F[\CBPart{F}{G}]]
        \subseteq A\setminus\CBPart{F}{G}}{%
      \FormulaRefAuto{CantorBernsteinComplementForward}[thm]{1,2}}
    \proofstepwidestar[1,2]{%
      A\setminus\CBPart{F}{G}
        \subseteq G[B]\setminus G[F[\CBPart{F}{G}]]}{%
      \FormulaRefAuto{CantorBernsteinComplementBackward}[thm]{1,2}}
    \proofstepwidestar[1,2]{%
      G[B]\setminus G[F[\CBPart{F}{G}]]
        =A\setminus\CBPart{F}{G}}{%
      \FormulaRefAuto{A\subseteq B\dsep B\subseteq A\vdash A=B}[thm]{3,4}}
  \closeproofpartwide
\end{tabproofsplitwide}
}
\Needspace{14\baselineskip}

\CSBSecondBranchImage
\CSBProof{%
\begin{tabproofwide}
  \proofstepwidestar[1]{F\colon A\inj B}{\rA}
  \proofstepwidestar[2]{G\colon B\inj A}{\rA}
  \proofstepwidestar[1,2]{\CBPart{F}{G}\subseteq A}{%
    \FormulaRefAuto{CantorBernsteinPartSubset}[thm]{1,2}}
  \proofstepwidestar[1]{F\colon A\to B}{%
    \FormulaRefAuto{Injektive Funktion}[def]{1}}
  \proofstepwidestar[1,2]{F[\CBPart{F}{G}]\subseteq B}{%
    \FormulaRefAuto{C\subseteq A\dsep F\colon A\to B
      \vdash F[C]\subseteq B}[thm]{3,4}}
  \proofstepwidestar[1,2]{B \subseteq B}{%
      \FormulaRefAuto{A\subseteq A}[thm]}
  \proofstepwidestar[1,2]{%
    G[B\setminus F[\CBPart{F}{G}]]
      =G[B]\setminus G[F[\CBPart{F}{G}]]}{%
    \FormulaRefAuto{InjectiveImageDifference}[thm]{2,5,
      6}}
  \proofstepwidestar[1,2]{%
    G[B]\setminus G[F[\CBPart{F}{G}]]
      =A\setminus\CBPart{F}{G}}{%
    \FormulaRefAuto{CantorBernsteinComplementIdentity}[thm]{1,2}}
  \proofstepwidestar[1,2]{%
    G[B\setminus F[\CBPart{F}{G}]]
      =A\setminus\CBPart{F}{G}}{\rChain{7,8}}
\end{tabproofwide}
}
\Needspace{14\baselineskip}

\CSBFixedInjections
\CSBProof{%
\begin{tabproofsplitwide}
\Needspace{8\baselineskip}
\proofpartwide{Erster bijektiver Zweig}
\proofstepwidestar[1]{F\colon A\inj B}{\rA}
  \proofstepwidestar[2]{G\colon B\inj A}{\rA}
  \proofstepwidestar[1,2]{\CBPart{F}{G}\subseteq A}{%
    \FormulaRefAuto{CantorBernsteinPartSubset}[thm]{1,2}}
  \proofstepwidestar[1,2]{%
    F\restriction_{\CBPart{F}{G}}\colon
      \CBPart{F}{G}\bij F[\CBPart{F}{G}]}{%
    \FormulaRefAuto{InjectionRestrictionToImageBijective}[thm]{1,3}}

\closeproofpartwide

\Needspace{8\baselineskip}

\proofpartwide{Zweiter bijektiver Zweig}
  \setcounter{proofstepnr}{4}% Fortlaufende Schritte dieses Beweises.
\proofstepwidestar[1]{F\colon A\to B}{%
    \FormulaRefAuto{Injektive Funktion}[def]{1}}
  \proofstepwidestar[1,2]{F[\CBPart{F}{G}]\subseteq B}{%
    \FormulaRefAuto{C\subseteq A\dsep F\colon A\to B
      \vdash F[C]\subseteq B}[thm]{3,5}}
  \proofstepwidestar[1,2]{B\setminus F[\CBPart{F}{G}]\subseteq B}{%
    \FormulaRefAuto{A\setminus B\subseteq A}[thm]}
  \proofstepwidestar[1,2]{%
    G\restriction_{B\setminus F[\CBPart{F}{G}]}\colon
      B\setminus F[\CBPart{F}{G}]
      \bij G[B\setminus F[\CBPart{F}{G}]]}{%
    \FormulaRefAuto{InjectionRestrictionToImageBijective}[thm]{2,7}}
  \proofstepwidestar[1,2]{%
    G[B\setminus F[\CBPart{F}{G}]]
      =A\setminus\CBPart{F}{G}}{%
    \FormulaRefAuto{CantorBernsteinSecondBranchImage}[thm]{1,2}}
  \proofstepwidestar[1,2]{%
    G\restriction_{B\setminus F[\CBPart{F}{G}]}\colon
      B\setminus F[\CBPart{F}{G}]
      \bij A\setminus\CBPart{F}{G}}{\rIE{9,8}}
  \proofstepwidestar[1,2]{%
    \bigl(G\restriction_{B\setminus F[\CBPart{F}{G}]}\bigr)^{-1}
      \colon A\setminus\CBPart{F}{G}
      \bij B\setminus F[\CBPart{F}{G}]}{%
    \FormulaRefAuto{InverseFunctionBijection}[thm]{10}}

\closeproofpartwide

\Needspace{8\baselineskip}

\proofpartwide{Disjunktheit und Verkleben}
  \setcounter{proofstepnr}{11}% Fortlaufende Schritte dieses Beweises.
\proofstepwidestar[]{A\setminus\CBPart{F}{G}\subseteq A\setminus\CBPart{F}{G}}{%
      \FormulaRefAuto{A\subseteq A}[thm]}
\proofstepwidestar[]{%
    \CBPart{F}{G}\cap(A\setminus\CBPart{F}{G})=\varnothing}{%
    \FormulaRefAuto{K\subseteq H\setminus M\vdash M\cap K=\varnothing}[thm]{%
      12}}
  \proofstepwidestar[]{B\setminus F[\CBPart{F}{G}]\subseteq B\setminus F[\CBPart{F}{G}]}{%
      \FormulaRefAuto{A\subseteq A}[thm]}
  \proofstepwidestar[]{%
    F[\CBPart{F}{G}]\cap
      (B\setminus F[\CBPart{F}{G}])=\varnothing}{%
    \FormulaRefAuto{K\subseteq H\setminus M\vdash M\cap K=\varnothing}[thm]{%
      14}}
  \proofstepwidestar[1,2]{%
    \CaseFun{\CBPart{F}{G}}{A\setminus\CBPart{F}{G}}
      {F\restriction_{\CBPart{F}{G}}}
      {\bigl(G\restriction_{B\setminus F[\CBPart{F}{G}]}\bigr)^{-1}}
      \colon
      \CBPart{F}{G}\cup(A\setminus\CBPart{F}{G})
      \bij F[\CBPart{F}{G}]
        \cup(B\setminus F[\CBPart{F}{G}])}{%
    \FormulaRefAuto{CaseFunBijectiveDisjointTargets}[thm]{13,15,4,11}}

\closeproofpartwide

\Needspace{8\baselineskip}

\proofpartwide{Identifikation von Definitions- und Zielbereich}
  \setcounter{proofstepnr}{16}% Fortlaufende Schritte dieses Beweises.
\proofstepwidestar[1,2]{%
    A=\CBPart{F}{G}\cup(A\setminus\CBPart{F}{G})}{%
    \FormulaRefAuto{M\subseteq H\vdash H=M\cup(H\setminus M)}[thm]{3}}
  \proofstepwidestar[1,2]{%
    B=F[\CBPart{F}{G}]\cup(B\setminus F[\CBPart{F}{G}])}{%
    \FormulaRefAuto{M\subseteq H\vdash H=M\cup(H\setminus M)}[thm]{6}}
  \proofstepwidestar[1,2]{%
    \CaseFun{\CBPart{F}{G}}{A\setminus\CBPart{F}{G}}
      {F\restriction_{\CBPart{F}{G}}}
      {\bigl(G\restriction_{B\setminus F[\CBPart{F}{G}]}\bigr)^{-1}}
      \colon A\bij B}{\rIE{17,\rIE{18,16}}}
  \proofstepwidestar[1,2]{\exists H\colon A\bij B}{\rEI{19}}

\closeproofpartwide
\end{tabproofsplitwide}
}
